% Abstract -- authored at integration. Nature-accessible + measured product pitch.
% British spelling, group "we", no superlatives, no em dashes. ~270 words.
Financial econometrics has a reproducibility problem that is, for the most part,
a software problem. Published results depend on undocumented data revisions, on
estimator details buried in analysis scripts, and on machine-specific numerics
that move the third decimal place. We describe \engine{}, a research programme
organised around a single compute engine that treats these as first-class
engineering concerns. The engine exposes a parameterised-only registry of
twenty-six methods spanning long-memory estimation, wavelet decomposition,
transfer entropy, surrogate testing, community detection, network connectedness,
and structural channel attribution. Every method ships with a pre-registered,
failable test of its own output, and a method is not in the catalogue until it
does. A nightly evaluation suite runs the methods on real market data against
fixed numerical bands, maintains a machine-written state ledger, and refreshes a
public claims record without ever editing a result by hand.

The engine runs across three planes: a scale-to-zero kernel on managed cloud
infrastructure, an always-on workstation tower that carries the heavy
asynchronous estimators, and static public surfaces that render only what the
suite has verified. We show that a governed core budget with pinned
linear-algebra threading makes results independent of worker count, and that a
GPU port of the $k$-nearest-neighbour transfer-entropy estimator reproduces the
CPU value to within machine epsilon while running several times faster. Six
published frameworks, on adaptive market networks, scale-ordered contagion,
contagion channels, wavelet quantile transfer, commodity systems, and
protocol-mediated systemic risk, are elements of this one programme rather than
separate codebases. We report honest negatives alongside positive findings, and
we argue that the engineering discipline is itself a contribution.
